%=======================================================================
\chapter{Kinematic Constraints}
%=======================================================================

Not every motion is available to every material. Water resists a change
of volume so much more stiffly than it resists a change of shape that,
over the range of pressures met in ordinary problems, it is both more
accurate and far more convenient to forbid the volume change outright
than to model it. A stiff fibre embedded in a soft matrix tells the same
story along a line rather than in bulk: the fibre will stretch, but by so
little that setting its stretch to one costs less error than any elastic
law for it would. Idealising a very stiff response as \emph{no} response
is what this chapter is about.

The idealisation is not free, and the price is the whole content of the
chapter. If the material cannot perform a certain motion, then something
must stop it, and that something is a stress. It is a stress of a new
kind: no constitutive equation produces it, because a constitutive
equation answers ``what stress does this deformation cause'', and here
the deformation in question never happens. What determines it instead is
the equations of motion together with the boundary and initial data. The
pressure in an incompressible fluid is the case everyone meets first, and
it is worth saying at the outset what will be proved: that pressure is
not a material property at all.

Guided by the incompressible fluid, we consider \emph{local} constraints
of the form
\begin{equation}
   g_{\kappa}(\bF)=0 ,
   \tag{10.1}\label{eq:10.1}
\end{equation}
from the perspective of an observer $O$, where $g_{\kappa}$ is a
scalar-valued function, $\bF$ is the deformation gradient of
Section~\ref{sec:defgrad}, and the subscript $\kappa$ records the
reference configuration relative to which $\bF$, and hence $g_{\kappa}$,
is defined. For incompressibility, $g_{\kappa}(\bF)=\det\bF-1$.

\begin{remark}[what a constraint is, and what it is not]\label{rem:ch10what}
Three distinctions are worth fixing before anything is computed, because
each of them is a way of misreading \eqr{10.1}.

\emph{A constraint is not a boundary condition.} A boundary condition
restricts the motion of particular material points, those on
$\partial\kappa$, and it is data supplied with the problem. Equation
\eqr{10.1} is imposed at \emph{every} $\bX\in\kappa$ and at every $t$,
and it is part of the description of the material, not of the problem.
Two different problems for the same incompressible fluid have different
boundary conditions and the same constraint.

\emph{A constraint is not a constitutive equation.} A constitutive
equation such as \eqr{8.33} is a map from the deformation to the stress;
it answers a question. A constraint answers nothing; it deletes part of
the domain on which the question may be asked. This is why the two must
be stated separately, and why \eqr{10.14} below will split the stress in
two rather than modifying $\hat\bT_{\kappa}$.

\emph{The constraint is on $\bF$, not on $\bchi_{\kappa}$.} The
deformation gradient is the local, first-order content of the motion at
a point, so \eqr{10.1} restricts what the material may do in an
arbitrarily small neighbourhood of $\bX$, which is what a statement about
material response has to do. A restriction on $\bchi_{\kappa}$ itself
would say where the particle is allowed to be, which is a statement about
the problem and not about the material. The same locality argument opened
Chapter 8.
\end{remark}

%=======================================================================
\section{Observer consensus}\label{sec:ch10obs}
%=======================================================================

Chapter 7 required all observers to agree on material response. Because a
constraint characterises the nature of that response, we require that all
observers agree that a constraint of the form \eqr{10.1} is in force.
Thus the constraint, as perceived by an observer $O^{+}$, is
\begin{equation}
   g^{+}_{\kappa^{+}}(\bF^{+})=0 ,
   \tag{10.2}\label{eq:10.2}
\end{equation}
where
\begin{equation}
   g^{+}_{\kappa^{+}}(\bF^{+})=g_{\kappa}(\bF).
   \tag{10.3}\label{eq:10.3}
\end{equation}
Equation \eqr{10.3} is the substantive assumption and \eqr{10.2} is its
consequence. The two observers are not merely required to reach the same
verdict, zero or non-zero; they are required to compute the same
\emph{number}. That is what ``agreeing on material response'' means when
the response is described by a scalar function.

From \eqr{7.17} the two deformation gradients are related by
$\bF^{+}=\bQ(t)\bF\bK^{t}$, in which
$\bQ(t)\colon\Vt\to\mathbb{E}^{3+}$ is the orthogonal tensor of the
change of observer and $\bK\colon\hVt\to\hat{\mathbb{E}}^{3+}$ is the
fixed orthogonal tensor of \eqr{7.15} relating the two reference
configurations. Solving for $\bF$,
\[
   \bQ^{t}\bF^{+}\bK=\bQ^{t}\bQ\bF\bK^{t}\bK=\I\bF\hI=\bF ,
\]
where $\bQ^{t}\bQ=\I$ is the identity on $\Vt$ and $\bK^{t}\bK=\hI$ the
identity on $\hVt$, both by orthogonality (Proposition~\ref{prop:orth}).
Substituting into \eqr{10.3},
\begin{equation}
   g^{+}_{\kappa^{+}}(\bF^{+})=g_{\kappa}(\bQ^{t}\bF^{+}\bK).
   \tag{10.4}\label{eq:10.4}
\end{equation}

As in Chapter 9, we now impose this requirement for \emph{two} relative
motions of the observers, both viewed from $O^{+}$. This is the device
that converts a statement relating two observers into a restriction on
the single function $g_{\kappa}$ used by one of them. Let the two
relative motions be characterised by orthogonal tensors $\bQ_{1}(t)$ and
$\bQ_{2}(t)$ in the sense of \eqr{9.6}. Observer $O^{+}$ perceives one
and the same $\bF^{+}$, so \eqr{10.4} applies twice with the same left
side and gives
\begin{equation}
   g_{\kappa}(\bQ_{2}^{t}\bF^{+}\bK)=g^{+}_{\kappa^{+}}(\bF^{+})
   =g_{\kappa}(\bQ_{1}^{t}\bF^{+}\bK).
   \tag{10.5}\label{eq:10.5}
\end{equation}
Here we assume $\bQ_{2}(t_{0})=\bQ_{1}(t_{0})$, so that $\bK$ is the same
for both relative motions; by \eqr{7.15}$_{1}$, $\bK$ is built from
$\bQ$ at the single instant $t_{0}$ and from the two shifters, so it is
this equality at $t_{0}$, and nothing else, that is needed. Let
$\bQ=\bQ_{2}^{t}\bQ_{1}\colon\Vt\to\Vt$, exactly as in \eqr{9.11}, and
let
\begin{equation}
   \bF=\bQ_{1}^{t}\bF^{+}\bK .
   \tag{10.6}\label{eq:10.6}
\end{equation}
Then, using $\bQ_{1}\bQ_{1}^{t}=\I^{+}$,
\[
   \bQ\bF=\bQ_{2}^{t}\bQ_{1}\bQ_{1}^{t}\bF^{+}\bK
        =\bQ_{2}^{t}\I^{+}\bF^{+}\bK=\bQ_{2}^{t}\bF^{+}\bK ,
\]
and taking determinants of the two ends, with
$\det(\bA\bB)=\det\bA\det\bB$ and $\det\bQ_{2}^{t}=\det\bQ_{2}$,
\begin{equation}
   \bQ_{2}^{t}\bF^{+}\bK=\bQ\bF;\qquad\text{hence}\qquad
   (\det\bQ)J=(\det\bQ_{2})(\det\bK)J^{+} ,
   \tag{10.7}\label{eq:10.7}
\end{equation}
in which $J=\det\bF$ and $J^{+}=\det\bF^{+}$.

Recall the discussion leading to \eqr{7.19}, on the choice of occupiable
reference configurations. Both $\bQ_{2}$ and $\bK$ are orthogonal, so
each determinant is $\pm1$ and their product is $\pm1$; occupiability
forces $J>0$ and $J^{+}>0$, and $J^{+}=J(\det\bQ_{2})(\det\bK)$ then
leaves only the choice $(\det\bQ_{2})(\det\bK)=1$, whence $J^{+}=J$.
Putting both into \eqr{10.7}$_{2}$ gives $(\det\bQ)J=J$, and since $J>0$
we may divide by it:
\[
   \det\bQ=1 .
\]
So $\bQ$ is a rotation. With \eqr{10.6} and \eqr{10.7}$_{1}$, equation
\eqr{10.5} now reads
\begin{equation}
   g_{\kappa}(\bQ\bF)=g_{\kappa}(\bF)
   \tag{10.8}\label{eq:10.8}
\end{equation}
for arbitrary rotations $\bQ$. The two observers have been eliminated
entirely: \eqr{10.8} is a statement about the one function $g_{\kappa}$
that $O$ uses.

\subsection*{Reduction of \eqr{10.8}}

To obtain a \emph{necessary} condition we exploit the arbitrariness of
$\bQ$ by choosing it as awkwardly as possible for the left side of
\eqr{10.8}. Take
\[
   \bQ=\shift\bR^{t},
\]
where $\shift$ is the shifter of Definition~\ref{def:shifter} and $\bR$
is the two-point rotation in the polar factorisation \eqr{3.32},
$\bF=\bR\bU$, of Theorem~\ref{thm:polar}. This is a legitimate choice.
Since $\bR\colon\hVt\to\Vt$ we have $\bR^{t}\colon\Vt\to\hVt$, and
$\shift\colon\hVt\to\Vt$, so the product maps $\Vt$ to $\Vt$ as $\bQ$
must; it is orthogonal as a product of orthogonal tensors, and its
determinant is $(\det\shift)(\det\bR)=1$ because both factors are proper.
With it,
\[
   \bQ\bF=\shift\bR^{t}\bR\bU=\shift\hI\bU=\shift\bU ,
\]
the middle step by $\bR^{t}\bR=\hI$. So \eqr{10.8} furnishes
\[
   g_{\kappa}(\bF)=g_{\kappa}(\shift\hI\bU)=g_{\kappa}(\shift\bU).
\]
Because $\shift$ is fixed in the frame of $O$, the right side is a
function of $\bU$ alone. The right stretch tensor is in turn uniquely
determined by the right Cauchy--Green tensor $\bC=\bF^{t}\bF$ through
\eqr{3.38}, $\bU=\sqrt{\bC}$, the square root being unique by
Theorem~\ref{thm:sqrt}. We therefore conclude that
\begin{equation}
   g_{\kappa}(\bF)=\bar g_{\kappa}(\bC)
   \tag{10.9}\label{eq:10.9}
\end{equation}
for some function $\bar g_{\kappa}$.

Conversely, suppose \eqr{10.9} holds. Then for every rotation $\bQ$,
\[
   g_{\kappa}(\bQ\bF)=\bar g_{\kappa}\big((\bQ\bF)^{t}\bQ\bF\big)
   =\bar g_{\kappa}(\bF^{t}\bQ^{t}\bQ\bF)
   =\bar g_{\kappa}(\bF^{t}\I\bF)=\bar g_{\kappa}(\bF^{t}\bF)
   =g_{\kappa}(\bF),
\]
so \eqr{10.9} is both necessary and sufficient for \eqr{10.8}.

\begin{remark}[what \eqr{10.9} forbids]\label{rem:ch10rotblind}
The content of \eqr{10.9} is negative and it is worth stating that way: a
kinematic constraint may not see the rotation $\bR$. Suppose someone
proposed the constraint $R_{11}=1$, that a certain material line never
turn out of a certain plane. Chapter 3 showed that $\bR$ is the part of
$\bF$ that carries no stretch at all, and Chapter 7 showed, in
\eqr{7.29}$_{3}$, that $\bR^{+}=\bQ\bR\bK^{t}$: the rotation is precisely
the part of $\bF$ that a change of observer alters. So the proposed
constraint would be satisfied for one observer and violated for another,
watching the same body. It would make ``this material is inextensible''
a statement about who is looking. Equation \eqr{10.9} is the exact
statement of what survives that objection: only $\bC$, and by
Proposition~\ref{prop:C} and Remark~\ref{rem:CBintuition} $\bC$ is the
part of the deformation that measures lengths and angles of material
fibres, which is what two observers can compare.
\end{remark}

The constraint from the perspective of $O^{+}$ now follows from
\eqr{10.4}. Here $\bQ\colon\Vt\to\mathbb{E}^{3+}$ is once more the
orthogonal tensor of the change of observer, and $\bK$ is a fixed
parameter:
\begin{equation}
\begin{aligned}
   g^{+}_{\kappa^{+}}(\bF^{+})
   &=g_{\kappa}(\bQ^{t}\bF^{+}\bK)
    =\bar g_{\kappa}\big((\bQ^{t}\bF^{+}\bK)^{t}\bQ^{t}\bF^{+}\bK\big)\\
   &=\bar g_{\kappa}\big(\bK^{t}(\bF^{+})^{t}\bQ\bQ^{t}\bF^{+}\bK\big)
    =\bar g_{\kappa}\big(\bK^{t}(\bF^{+})^{t}\I^{+}\bF^{+}\bK\big)\\
   &=\bar g_{\kappa}\big(\bK^{t}(\bF^{+})^{t}\bF^{+}\bK\big)
    =\bar g_{\kappa}(\bK^{t}\bC^{+}\bK)
    =\bar g^{+}_{\kappa^{+}}(\bC^{+}) .
\end{aligned}
   \tag{10.10}\label{eq:10.10}
\end{equation}
The transpose of a product reverses the order, which produced the second
line; $\bQ\bQ^{t}=\I^{+}$ is orthogonality read in the other direction
from the one used above, and it is the step at which the observer
disappears. What remains, $\bK^{t}\bC^{+}\bK$, involves only the two
reference configurations, and $\bK$ does not depend on time. Comparing
with \eqr{7.30}$_{1}$, $\bC^{+}=\bK\bC\bK^{t}$, the last line simply says
$\bar g^{+}_{\kappa^{+}}(\bC^{+})=\bar g_{\kappa}(\bC)$, which is
\eqr{10.3} again with $\bC$ in place of $\bF$.

\subsection*{Examples}

\subsubsection*{1. Incompressibility}

Because $J>0$, the constraint $J=1$ associated with incompressibility is
equivalent to $1=J^{2}$, and
\[
   J^{2}=(\det\bF)^{2}=(\det\bF^{t})(\det\bF)=\det(\bF^{t}\bF)=\det\bC ,
\]
using $\det\bA^{t}=\det\bA$ and the product rule for determinants.
Hence
\begin{equation}
   \bar g_{\kappa}(\bC)=\det\bC-1 .
   \tag{10.11}\label{eq:10.11}
\end{equation}
The step $J=1\iff J^{2}=1$ needs $J>0$ and would be false without it;
$J>0$ is guaranteed for a motion by Remark~\ref{rem:vol} and the
requirement that material not turn itself inside out. Note also that
\eqr{10.11} is manifestly of the form \eqr{10.9}, as it had to be:
$\det\bF$ by itself is \emph{not} an admissible constraint function in
the sense of \eqr{10.9}, but on rotations $\det\bQ=1$ so
$\det(\bQ\bF)=\det\bF$ and \eqr{10.8} holds anyway. The two statements
agree, as the equivalence just proved says they must.

\subsubsection*{2. Inextensibility}

In the engineering theory of fibre-reinforced composites the fibres are
idealised as continuously distributed material curves in $\kappa$ with a
unit tangent field $\bM(\bX)$. If the fibres are very stiff compared with
the matrix in which they are embedded, their stretches $\mu$ stay close
to unity in any deformation the composite can survive, and it is then
appropriate to impose $\mu=1$ exactly. By \eqr{2.83}, $\mu^{2}
=\ip\bM{\bC\bM}$, so with $\mu>0$ the constraint $\mu=1$ is
$\ip\bM{\bC\bM}=1$ and
\begin{equation}
   \bar g_{\kappa}(\bC)=\ip\bM{\bC\bM}-1 .
   \tag{10.12}\label{eq:10.12}
\end{equation}
Here $\bM$ is a material vector in the sense of
Definition~\ref{def:matvec}: it labels a fibre, and by
Remark~\ref{rem:matvec} it satisfies $\dot\bM=\bze$. That fact will be
used when \eqr{10.12} is differentiated.

\subsubsection*{3. Rigidity}

If the body is rigid then $\bF$ is a spatially uniform rotation
(Section~\ref{sec:rigid}), so $\bC=\bF^{t}\bF=\hI$. Because $\bC$ is
symmetric, this single tensor equation is equivalent to six scalar
constraints,
\begin{equation}
   \bar g_{(\kappa)AB}=C_{AB}-\dd_{AB},
   \tag{10.13}\label{eq:10.13}
\end{equation}
one for each of the six independent entries of the symmetric matrix
$(C_{AB})$.

\begin{remark}[six is the maximum]\label{rem:ch10sixmax}
There can be no more than six independent constraints, because there are
at most six independent entries in $(C_{AB})$. Rigidity is therefore the
\emph{maximal} constraint: it fixes every component of that matrix and
leaves the material no freedom whatever beyond a rigid motion. The number
six is the dimension of $\Sym$, computed in Chapter 1 by counting
$C_{11},C_{22},C_{33},C_{12},C_{13},C_{23}$, and the same six will
reappear in \eqr{10.19} as the maximum number of multipliers. That is not
a coincidence and Section~\ref{sec:ch10stress} explains why: each
independent constraint deletes one dimension from the space of admissible
stretchings and pays for it with one dimension of reactive stress, and
both spaces sit inside the same six-dimensional $\Sym$.
\end{remark}

%=======================================================================
\section{Modification of the stress to accommodate
constraints}\label{sec:ch10stress}
%=======================================================================

We have seen, in connection with the incompressible Newtonian fluid of
Chapter 9, that the stress decomposes into the sum of a reactive part
$-p\I$ and a constitutive part determined by the material. The salient
feature of the reactive part is that it makes no contribution to the
stress power density whenever the constraint is operative: by
Definition~\ref{def:tip}, $-p\I\cdot\bD=-p\tr\bD$, which vanishes
whenever $\tr\bD=0$. We extend this idea to general constraints by
assuming that the net Cauchy stress is
\begin{equation}
   \bT=\bN+\bG ,
   \tag{10.14}\label{eq:10.14}
\end{equation}
where $\bG$ is a symmetric-tensor-valued constitutive function
appropriate to the material and the conditions at hand, and $\bN$ is a
symmetric tensor such that $\bN\cdot\bD$ vanishes in any motion
compatible with the constraint. Once the structure of $\bN$ has been
determined, any indeterminacy remaining in it must be such as to secure a
solution, if one exists, to the equations of motion together with the
boundary and initial conditions.

\begin{remark}[why worklessness, and why it is an assumption]
\label{rem:ch10workless}
Two questions press here. Why should the reaction do no work, and is that
a theorem or a hypothesis?

\emph{Why worklessness is the right requirement.} The stress power of a
part $S$ is $\mathcal{S}(S,t)=\int_{P_{t}}\bT\cdot\bL\,\dl v$ by
\eqr{6.153}. Since $\bT$ is symmetric and $\bW$ is skew,
$\bT\cdot\bW=0$ by Lemma~\ref{lem:symskw}, so with $\bL=\bD+\bW$ from
\eqr{4.7} the stress power density is $\bT\cdot\bL=\bT\cdot\bD$. Under
\eqr{10.14} this is $\bN\cdot\bD+\bG\cdot\bD$. If $\bN\cdot\bD$ vanished
only sometimes, the reaction would be feeding energy into or draining it
from the body, and by the energy balance \eqr{6.168} that energy would
have to show up as heat or internal energy. A constraint is an
idealisation of a very stiff \emph{elastic} response, which stores energy
and gives it back and dissipates none; in the limit the stored energy
goes to zero along with the deformation that stores it. Requiring
$\bN\cdot\bD=0$ is the exact expression of that limit.

The same conclusion arrives from the thermodynamic side. The
Clausius--Duhem inequality carries the stress only through $\bT\cdot\bL$,
so $\bN$ drops out of it identically. Thermodynamics therefore places no
restriction whatever on $\bN$, which is precisely the sentence in Chapter
9 that the spherical part of the stress in an incompressible fluid ``is
constitutively indeterminate and hence outside the purview of our earlier
discussion about observer consensus concerning material response''.
Worklessness and constitutive indeterminacy are two readings of one fact.

\emph{It is an assumption.} Nothing proved so far forces \eqr{10.14} on
us; a theory in which the constraint reaction did work could be written
down and would not be self-contradictory. It would, however, require
extra constitutive information to say how much work, and the whole
economy of the constraint idea is that it requires none. This is the
continuum version of the principle of d'Alembert, that the forces of
constraint in a mechanical system do no virtual work, and it is adopted
for the same reason: it is the assumption that makes constraints cost
nothing to describe.
\end{remark}

Suppose, then, that the constraint is operative at a particular material
point throughout an interval $I$ of time. We require that
\begin{equation}
   \bN\cdot\bD=0\quad\text{for all }\bD\text{ such that }
   \bar g_{\kappa}(\bC(t))=0;\qquad t\in I .
   \tag{10.15}\label{eq:10.15}
\end{equation}

Because $\bar g_{\kappa}$ vanishes identically on $I$, the function
$t\mapsto\bar g_{\kappa}(\bC(t))$ is the zero function there, and the
derivative of the zero function is zero. Writing $g$ in place of
$\bar g_{\kappa}$ to keep the notation tractable, and using the chain
rule for a scalar-valued function of a tensor,
\begin{equation}
   0=\dot g=g_{\bC}\cdot\dot\bC=\tr\big((g_{\bC})\dot\bC\big),
   \tag{10.16}\label{eq:10.16}
\end{equation}
where $g_{\bC}=(g_{\bC})^{t}$ is the symmetric referential tensor-valued
gradient of $g$ (Sections B.2 and B.3 of Appendix B), defined by
$\dl g=g_{\bC}\cdot\dl\bC$ and existing uniquely by the Riesz argument of
Theorem~\ref{thm:riesz} applied on the inner-product space $\Lin$.

Three points in \eqr{10.16} need spelling out.

\emph{Why the interval, and not the instant.} The constraint must hold
throughout $I$ for the differentiation to be legal. Had \eqr{10.15}
required only $\bar g_{\kappa}(\bC(t_{0}))=0$ at one instant, the
function $t\mapsto\bar g_{\kappa}(\bC(t))$ would be zero at a single
point and its derivative there would be unrestricted. A constraint is by
its nature a standing condition, so this costs nothing, but it is the
hypothesis the argument actually uses.

\emph{Why $g_{\bC}$ may be taken symmetric.} By
Proposition~\ref{prop:C}, $\bC(t)$ is symmetric at every $t$, so each
difference quotient $\big(\bC(t+h)-\bC(t)\big)/h$ is symmetric, and
$\Sym$ is a subspace and therefore closed: the limit $\dot\bC$ is
symmetric. By the orthogonality of the splitting
$\Lin=\Sym\oplus\Skw$ of Definition~\ref{def:splits} and
Lemma~\ref{lem:symskw}, $\Skw(g_{\bC})\cdot\dot\bC=0$, so the inner
product in \eqr{10.16} picks out only $\Sym(g_{\bC})$. The skew part of
$g_{\bC}$ is invisible to \eqr{10.16}, which is what effectively defines
$g_{\bC}$; no generality is lost by taking it symmetric, and doing so
makes it unique.

\emph{The order of operations.} The function $g(\bC)$ is well defined for
\emph{all} symmetric positive definite $\bC$, not only for those that
satisfy the constraint, and $g_{\bC}$ is therefore likewise well defined
for all such $\bC$. The derivative appearing in \eqr{10.16} is obtained
by differentiating first and imposing the constraint afterwards.

\begin{remark}[what breaks if the constraint is imposed too
early]\label{rem:ch10apost}
Reverse the order and the entire theory collapses in one line. Impose the
constraint first: then $g$ is the zero function on the set where it is
defined, its gradient is $\hOb$, and \eqr{10.18} below gives $\bN=\Ob$.
There would be no reaction, the incompressible fluid would carry no
pressure, and a rigid body would be free of stress.

The error is the same one that Remark~\ref{rem:whykt0} of Chapter 4
warned against in another guise: confusing the value of a function on a
set with the function itself. What the constraint says is that the motion
never leaves the surface $g=0$ in the space of symmetric positive
definite tensors. The gradient $g_{\bC}$ is the \emph{normal} to that
surface, and a normal is computed from how $g$ behaves as one steps
\emph{off} the surface, which is exactly the information deleted by
imposing the constraint first. That is why \eqr{10.21} is computed for
all $\bC$ and only then evaluated, in \eqr{10.22}, on the constraint set.
\end{remark}

Now bring in Chapter 4. Recalling \eqr{4.11}, $\bF^{t}\bD\bF=\tfrac12
\dot\bC$, that is $\dot\bC=2\bF^{t}\bD\bF$, equation \eqr{10.16} becomes
$0=2\tr\big(g_{\bC}(\bF^{t}\bD\bF)\big)$, and dividing by the constant
$2$,
\begin{equation}
   0=\tr\big(g_{\bC}(\bF^{t}\bD\bF)\big)
    =\tr\big((\bF(g_{\bC})\bF^{t})\bD\big)
    =\bF(g_{\bC})\bF^{t}\cdot\bD ,
   \tag{10.17}\label{eq:10.17}
\end{equation}
in which the first factor on the right-hand side of the third equality is
a symmetric \emph{spatial} tensor.

The middle equality of \eqr{10.17} is the one step here that is not
purely formal, since the two products live on different spaces and the
trace of a two-point tensor is not defined. In components, with
$g_{\bC}=G_{AB}\hbE_{A}\otimes\hbE_{B}$, $\bF=F_{iA}\be_{i}
\otimes\hbE_{A}$ and $\bD=D_{ij}\be_{i}\otimes\be_{j}$, the tensor
$g_{\bC}\bF^{t}\bD\bF$ is referential with components
$G_{AB}F_{iB}D_{ij}F_{jE}$, so its trace is
\[
   \tr\big(g_{\bC}\bF^{t}\bD\bF\big)=G_{AB}F_{iB}D_{ij}F_{jA},
\]
whereas $\bF(g_{\bC})\bF^{t}$ is spatial with components
$F_{iA}G_{AB}F_{jB}$, so
\[
   \bF(g_{\bC})\bF^{t}\cdot\bD=F_{iA}G_{AB}F_{jB}D_{ij}.
\]
The two are the same sum of products of the same numbers, written in a
different order. Both traces are legitimate, each on its own space, and
they are equal. Note also that
$\big(\bF(g_{\bC})\bF^{t}\big)^{t}=\bF(g_{\bC})^{t}\bF^{t}
=\bF(g_{\bC})\bF^{t}$, so the tensor is symmetric, and that the index
bookkeeping of Remark~\ref{rem:hatinch4} confirms it is spatial: $\bF$
carries $\be_{i}\otimes\hbE_{A}$, $g_{\bC}$ carries
$\hbE_{A}\otimes\hbE_{B}$, $\bF^{t}$ carries $\hbE_{B}\otimes\be_{j}$,
and the capital indices contract in pairs leaving
$\be_{i}\otimes\be_{j}$.

Equation \eqr{10.17} says exactly this: the stretchings $\bD$ that the
constraint permits are those, and only those, perpendicular in $\Sym$ to
the fixed symmetric spatial tensor $\bF(g_{\bC})\bF^{t}$. A constraint on
$\bC$, which is a statement about the deformation as a whole, has become
a single linear equation for $\bD$, which is a statement about the
present rate. The bridge is \eqr{4.11}, and this is the sense in which
the constraint restricts the stretching tensor of Chapter 4.

\begin{remark}[from a surface to a plane]\label{rem:ch10lin}
The set of $\bC$ satisfying $g(\bC)=0$ is a curved five-dimensional
surface in the six-dimensional $\Sym$; the set of $\bD$ satisfying
\eqr{10.17} is a flat five-dimensional subspace. Nothing was lost in
passing from one to the other, because a velocity is by definition
tangent to the surface it moves on, and the tangent space to a level
surface at a point is a subspace. Linearity is what makes the next step
possible at all: one may speak of the orthogonal complement of a
subspace, but not of the orthogonal complement of a curved surface.
\end{remark}

\subsection*{The reaction is pinned down to one multiplier}

By \eqr{10.15} we require $\bN$ to be perpendicular, in $\Sym$, to every
$\bD$ that is perpendicular to $\bF(g_{\bC})\bF^{t}$
(Figure~\ref{fig:ch10ortho}). Accordingly $\bN$ must be parallel to
$\bF(g_{\bC})\bF^{t}$. The following lemma is that inference, made
precise; it is the same manoeuvre as Lemma~\ref{lem:cvec} of Chapter 5,
namely, feed the hypothesis the one special element that turns it into a
statement about a norm.

\begin{lemma}[the double orthogonal complement of a line]
\label{lem:ch10ortho}
Let $\bA\in\Sym$ with $\bA\neq\Ob$, and let $\bN\in\Sym$ satisfy
$\bN\cdot\bD=0$ for every $\bD\in\Sym$ with $\bA\cdot\bD=0$. Then
$\bN=\lambda\bA$ for exactly one scalar $\lambda$.
\end{lemma}

\begin{proof}
By Theorem~\ref{thm:tip} the tensor inner product is positive definite,
so $\bA\cdot\bA=|\bA|^{2}>0$ and we may set
\[
   \lambda=\frac{\bN\cdot\bA}{\bA\cdot\bA},\qquad
   \bN'=\bN-\lambda\bA .
\]
Then $\bN'\in\Sym$, being a difference of symmetric tensors, and
\[
   \bA\cdot\bN'=\bA\cdot\bN-\lambda\,\bA\cdot\bA
   =\bA\cdot\bN-\frac{\bN\cdot\bA}{\bA\cdot\bA}\,\bA\cdot\bA=0 .
\]
So $\bN'$ is one of the $\bD$'s covered by the hypothesis, and taking
$\bD=\bN'$ gives $\bN\cdot\bN'=0$. On the other hand
$\bN=\bN'+\lambda\bA$, so
\[
   0=\bN\cdot\bN'=(\bN'+\lambda\bA)\cdot\bN'
   =|\bN'|^{2}+\lambda\,(\bA\cdot\bN')=|\bN'|^{2},
\]
whence $\bN'=\Ob$ by positive definiteness again, and $\bN=\lambda\bA$.
Uniqueness: if $\lambda_{1}\bA=\lambda_{2}\bA$ then
$(\lambda_{1}-\lambda_{2})\bA=\Ob$, and $\bA\neq\Ob$ forces
$\lambda_{1}=\lambda_{2}$.
\end{proof}

Applying the lemma with $\bA=\bF(g_{\bC})\bF^{t}$,
\begin{equation}
   \bN=\lambda(\bx,t)\,\bF(g_{\bC})\bF^{t},
   \tag{10.18}\label{eq:10.18}
\end{equation}
in which the multiplier $\lambda(\bx,t)$ is determined solely by the
equations of motion, with the stress given by \eqr{10.14}, together with
the constraint equation and the initial and boundary conditions.

\begin{remark}[the multiplier is a field, and it is a Lagrange
multiplier]\label{rem:ch10lagrange}
Read \eqr{10.18} beside the elementary rule for constrained
stationarity and it is the same statement. To make $f$ stationary on the
surface $g=0$ one sets $\nabla f=\lambda\nabla g$: the reaction is
directed along the normal to the constraint surface, and $\lambda$ is
whatever number the constraint then requires. Here $g_{\bC}$ is the
normal to the constraint surface in the referential space of symmetric
tensors, the sandwich $\bF(\,\cdot\,)\bF^{t}$ carries it to the spatial
one, and \eqr{10.18} says the reaction stress lies along it.

Two features of $\lambda$ are easy to misread. It is a \emph{field},
$\lambda(\bx,t)$, not a constant and not a material parameter: it is
allowed to vary from point to point and instant to instant, and it must,
since it is doing the work of enforcing the constraint everywhere at
once. And it is \emph{not} a function of $\bC$, $\bF$ or $\bD$. Writing
$\lambda=\hat\lambda(\bC)$ would smuggle a constitutive equation back in
and would generally make the resulting system unsolvable, for the reason
Chapter 9 gives after \eqr{9.85}: without a free field the equations for
an incompressible fluid outnumber the unknowns.
\end{remark}

If $n$ constraints $g^{(i)}(\bC)=0$, $i=1,\dots,n$, are in force during
an interval of time, then $\bN\cdot\bD$ must vanish for all $\bD$ such
that $\bF(g^{(i)}_{\bC})\bF^{t}\cdot\bD$ vanishes for each $i$. The $n$
constraints are said to be independent provided that the elements of
$\{\bF(g^{(i)}_{\bC})\bF^{t}\}$ are linearly independent, the
$g^{(i)}_{\bC}$ being computed for all positive definite symmetric $\bC$.
In this case $n\leq6$, because $\Sym$ is six-dimensional and
$\mathrm{Span}\{\bF(g^{(i)}_{\bC})\bF^{t}\}$ is then a subspace of
$\Sym$. Thus
\begin{equation}
   \bN=\sum_{i=1}^{n\leq6}\lambda_{i}(\bx,t)\,\bF(g^{(i)}_{\bC})\bF^{t}
   =\bF\Big(\sum_{i=1}^{n}\lambda_{i}g^{(i)}_{\bC}\Big)\bF^{t},
   \tag{10.19}\label{eq:10.19}
\end{equation}
where the multipliers $\lambda_{i}(\bx,t)$ are again such as to secure a
solution to the equations of motion together with the $n$ constraint
equations and the initial and boundary conditions, and the derivatives
$g^{(i)}_{\bC}$ are evaluated, after the fact, at those $\bC$'s that
satisfy all $n$ constraints simultaneously. The second form of
\eqr{10.19} follows from the first because $\bF(\,\cdot\,)\bF^{t}$ is
linear in its argument.

The proof of \eqr{10.19} is Lemma~\ref{lem:ch10ortho} with a line
replaced by a subspace. Let $\mathcal{W}=\mathrm{Span}
\{\bF(g^{(i)}_{\bC})\bF^{t}\}\subset\Sym$. The admissible $\bD$ are those
in $\mathcal{W}^{\perp}$, and the hypothesis is
$\bN\in(\mathcal{W}^{\perp})^{\perp}$. Decompose
$\bN=\bN_{\mathcal{W}}+\bN'$ with $\bN_{\mathcal{W}}\in\mathcal{W}$ and
$\bN'\in\mathcal{W}^{\perp}$, which is possible because $\mathcal{W}$ is
a finite-dimensional subspace of an inner-product space. Then $\bN'$ is
an admissible $\bD$, so $0=\bN\cdot\bN'=|\bN'|^{2}$ and $\bN'=\Ob$;
therefore $\bN\in\mathcal{W}$, which is \eqr{10.19}.

\bookfig{1}\begin{figure}[htbp]\centering
\begin{tikzpicture}[scale=1.0,
   ar/.style={-{Stealth[length=2.4mm]},thick}]
% the five-dimensional plane of admissible D, drawn as a parallelogram
\draw[thick,fill=blue!6]
   (-3.1,-0.95) -- (2.1,-0.95) -- (3.1,0.95) -- (-2.1,0.95) -- cycle;
\coordinate (O) at (0,0);
\fill (O) circle (1.3pt);
% normal direction, with its dashed continuation: the line of reactions
\draw[guide] (0,-1.75) -- (0,0);
\draw[ar] (O) -- (0,2.45)
   node[right,font=\small] {$\bF(g_{\bC})\bF^{t}$};
\draw[ar,red!65!black] (O) -- (0,1.35)
   node[left,font=\small,text=red!65!black]
   {$\bN=\lambda\bF(g_{\bC})\bF^{t}$};
% two admissible stretchings inside the plane
\draw[ar,blue!55!black] (O) -- (2.05,0)
   node[right=1pt,font=\small,text=blue!55!black] {$\bD$};
\draw[ar,blue!55!black] (O) -- (0.85,0.72)
   node[above right=-2pt,font=\small,text=blue!55!black] {$\bD'$};
% right angle mark between the vertical normal and the horizontal D
\draw[semithick] (0.26,0) -- (0.26,0.26) -- (0,0.26);
% labels
\node[font=\footnotesize,text=gray!85,align=left,anchor=west]
   at (-3.05,-1.45) {admissible stretchings:
   $\bF(g_{\bC})\bF^{t}\cdot\bD=0$};
\node[font=\footnotesize,text=gray!85,anchor=west] at (0.28,2.05)
   {reactions};
\node[font=\small] at (2.55,1.35) {$\Sym$};
\end{tikzpicture}
\caption{Orthogonality of $\bF(g_{\bC})\bF^{t}$ and $\bD$ in the space of
symmetric tensors. One constraint deletes one dimension from $\Sym$: the
admissible stretchings fill the five-dimensional plane through the origin
perpendicular to $\bF(g_{\bC})\bF^{t}$, drawn here as a parallelogram.
Requiring $\bN$ to be perpendicular to all of them leaves exactly the
one-dimensional line along $\bF(g_{\bC})\bF^{t}$, which is \eqr{10.18}.
The picture also shows why the count in Remark~\ref{rem:ch10sixmax} comes
out as it does: dimensions lost by $\bD$ are dimensions gained by $\bN$,
and $5+1=6=\dim\Sym$.}
\label{fig:ch10ortho}
\end{figure}

\subsection*{Examples}

We determine the forms of the reactive stresses for the constraints of
incompressibility, inextensibility and rigidity.

\subsubsection*{(1) Incompressibility}

In the case of incompressibility we use \eqr{10.11}, with \eqr{B.26} and
\eqr{B.27} of Appendix B, to obtain
\begin{equation}
   g_{\bC}\cdot\dot\bC=(\det\bC)^{\textstyle\cdot}
   =\bC^{*}\cdot\dot\bC=(\det\bC)\bC^{-1}\cdot\dot\bC ,
   \tag{10.20}\label{eq:10.20}
\end{equation}
for arbitrary symmetric $\dot\bC$.

Both of the equalities after the first deserve a derivation, since they
are the whole computation. Start from the box-product characterisation of
the determinant, Theorem~\ref{thm:det}: for any invertible $\bA(t)$ and
fixed $\bu,\bv,\bw$,
\[
   (\det\bA)\tbox\bu\bv\bw=\tbox{\bA\bu}{\bA\bv}{\bA\bw}.
\]
Differentiating, the box product being trilinear so that its derivative
obeys a three-term product rule, and writing $\ba=\bA\bu$, $\bb=\bA\bv$,
$\bc=\bA\bw$ together with $\dot\bA\bu=\dot\bA\bA^{-1}\ba$ and likewise
for the others,
\[
\begin{aligned}
   (\det\bA)^{\textstyle\cdot}\tbox\bu\bv\bw
   &=\tbox{\dot\bA\bu}{\bA\bv}{\bA\bw}+\tbox{\bA\bu}{\dot\bA\bv}{\bA\bw}
     +\tbox{\bA\bu}{\bA\bv}{\dot\bA\bw}\\
   &=\tbox{\dot\bA\bA^{-1}\ba}\bb\bc+\tbox\ba{\dot\bA\bA^{-1}\bb}\bc
     +\tbox\ba\bb{\dot\bA\bA^{-1}\bc}\\
   &=\tr\big(\dot\bA\bA^{-1}\big)\tbox\ba\bb\bc
    =\tr\big(\dot\bA\bA^{-1}\big)(\det\bA)\tbox\bu\bv\bw ,
\end{aligned}
\]
the third line by Definition~\ref{def:inv3}$_{1}$, which is precisely
that sum of three box products. Cancelling the non-zero
$\tbox\bu\bv\bw$,
\[
   (\det\bA)^{\textstyle\cdot}=(\det\bA)\tr\big(\dot\bA\bA^{-1}\big).
\]
This is Problem 2(a) of Chapter 4 with $\bA=\bF$; here we want it with
$\bA=\bC$. Using $\tr(\bA\bB)=\tr(\bB\bA)$ and then the symmetry of
$\bC^{-1}$ and $\dot\bC$,
\[
   (\det\bC)^{\textstyle\cdot}=(\det\bC)\tr\big(\dot\bC\bC^{-1}\big)
   =(\det\bC)\tr\big(\bC^{-1}\dot\bC\big)
   =(\det\bC)\,\bC^{-1}\cdot\dot\bC ,
\]
which is the last equality of \eqr{10.20}. The middle one is then
Proposition~\ref{prop:I2cof}, $\bA^{*}=(\det\bA)\bA^{-t}$, read for the
symmetric $\bC$: $\bC^{*}=(\det\bC)\bC^{-1}$. So the cofactor of
Definition~\ref{def:cof} is the derivative of the determinant, which is
the same fact that produced $\partial J/\partial F_{iA}=F^{*}_{iA}$ in
Chapter 4.

Since \eqr{10.20} holds for arbitrary symmetric $\dot\bC$ and both
$g_{\bC}$ and $(\det\bC)\bC^{-1}$ are symmetric, their difference is a
symmetric tensor orthogonal to every symmetric tensor, hence orthogonal
to itself, hence zero. Thus
\begin{equation}
   g_{\bC}=(\det\bC)\bC^{-1},
   \tag{10.21}\label{eq:10.21}
\end{equation}
which reduces to
\begin{equation}
   g_{\bC}=\bC^{-1}
   \tag{10.22}\label{eq:10.22}
\end{equation}
on the set of $\bC$'s that satisfy the constraint, where $\det\bC=1$.
Equations \eqr{10.21} and \eqr{10.22} are Remark~\ref{rem:ch10apost} in
action: the first is computed off the constraint set, the second is the
first evaluated on it, and doing those two steps in the other order would
have produced $g_{\bC}=\hOb$.

Writing $\lambda=-\bar p$, the associated reactive stress follows from
\eqr{10.18}:
\begin{equation}
   \bN=-\bar p\,\bF\bC^{-1}\bF^{t}=-\bar p\,\bF(\bF^{t}\bF)^{-1}\bF^{t}
   =-\bar p\,\I ,
   \tag{10.23}\label{eq:10.23}
\end{equation}
and is therefore mechanically equivalent to a pressure. The middle step
substituted $\bC=\bF^{t}\bF$; the last used
$(\bF^{t}\bF)^{-1}=\bF^{-1}\bF^{-t}$, from
Proposition~\ref{prop:prodinv}, so that
\[
   \bF(\bF^{t}\bF)^{-1}\bF^{t}=\bF\bF^{-1}\bF^{-t}\bF^{t}
   =\I\,\I=\I .
\]

The worklessness may now be verified directly rather than assumed. By
Definition~\ref{def:tip} and Problem 2(a) of Chapter 4,
\[
   \bN\cdot\bD=-\bar p\,\I\cdot\bD=-\bar p\tr\bD=-\bar p\dvg\bv
   =-\bar p\,\frac{\dot J}{J},
\]
which vanishes because the constraint $J\equiv1$ gives $\dot J=0$. The
reaction does no work for the plainest of reasons: a pressure works only
through a change of volume, and the constraint is that there is none.

For example, in the case of the Newtonian viscous fluid model \eqr{9.60},
whose stress is the sum of $-\tilde p\I$, a term proportional to
$(\tr\bD)\I$, and $2\mu\bD$, the constraint annihilates the middle term
and leaves the constitutive part $\bG=-\tilde p\I+2\mu\bD$. Adding the
reaction \eqr{10.23} and writing $p=\tilde p+\bar p$,
\begin{equation}
   \bT=-p\,\I+2\mu\bD
   \tag{10.24}\label{eq:10.24}
\end{equation}
if the constraint is operative. The constitutive indeterminacy of
$\bar p$ implies that $p$ is likewise indeterminate, and we arrive at
\eqr{9.83}. Note what has become of the thermodynamic pressure
$\tilde p=\rho^{2}\psi_{\rho}$ of \eqr{9.74}: it has not been
contradicted, it has been absorbed. Since $\bar p$ is arbitrary so is the
sum, and there is no way to recover $\tilde p$ from a measurement of
$\bT$.

\begin{example}[the pressure in a fluid at rest, and why it is not a
material property]\label{ex:ch10hydro}
This is the example that makes the whole chapter concrete, so it is
worked out completely.

\emph{The problem.} Let the fluid occupy $0\leq x_{3}\leq h$, be at rest,
and be acted on by gravity $\bb=-g\be_{3}$. At rest $\bv=\bze$, so
$\bL=\Ob$ and $\bD=\Ob$; the constraint $\tr\bD=0$ holds, and
\eqr{10.24} gives $\bT=-p\I$.

\emph{The equation of motion.} Substituting into
$\dvg\bT+\rho\bb=\rho\dot\bv$, with $\dot\bv=\bze$ because the fluid is
at rest. In components,
\[
   (\dvg\bT)_{i}=T_{ij,j}=(-p\dd_{ij})_{,j}=-p_{,i},
   \qquad\text{that is}\qquad \dvg(-p\I)=-\grad p ,
\]
so
\[
   \grad p=\rho\bb=-\rho g\be_{3},
   \qquad\text{i.e.}\qquad
   p_{,1}=0,\quad p_{,2}=0,\quad p_{,3}=-\rho g .
\]
The first two say $p$ depends on $x_{3}$ alone; integrating the third,
$p=c-\rho gx_{3}$ with $c$ a constant of integration.

\emph{The boundary condition.} On the free surface $x_{3}=h$ the exterior
unit normal is $\be_{3}$ and the air exerts the traction $-p_{a}\be_{3}$.
Cauchy's relation gives $\bt=\bT\be_{3}=-p\be_{3}$, so $p(h)=p_{a}$ and
$c=p_{a}+\rho gh$. Therefore
\[
   \boxed{\ p=p_{a}+\rho g(h-x_{3}).\ }
\]

\emph{The point.} Not one step of that calculation consulted the
material. The viscosity $\mu$ never appeared, because $\bD=\Ob$; the
density entered only through the body force, and $\tilde p$ was absorbed
into $p$ at the start. What fixed $p$ was the equation of motion, which
is a balance law of Chapter 6 and holds for every material, together with
one boundary condition.

To see that this is not an accident of the example, seal the container
with a piston and press on it, so that the boundary condition at
$x_{3}=h$ becomes $p=p_{a}+\Delta$ (Figure~\ref{fig:ch10press}). The
material is the same, the motion is the same, $\bD$ is still $\Ob$, and
$\bG$ is unchanged; yet now
\[
   p=p_{a}+\Delta+\rho g(h-x_{3}).
\]
Two different stress fields in the same material undergoing the same
motion. No constitutive equation can be responsible for the difference,
because a constitutive equation is a function of the motion and the
motion did not change. The pressure in an incompressible fluid is a
reaction, and $\Delta$ is the price of the constraint.

\emph{A dynamic case, so that the split is visible.} Take instead the
shear flow $\bv=\gamma x_{2}\be_{1}$ with $\gamma$ constant. Then
$\bL=\gamma\be_{1}\otimes\be_{2}$, so $\tr\bL=0$ and the flow is
isochoric as the constraint demands, and by \eqr{4.7}
$\bD=\tfrac12\gamma(\be_{1}\otimes\be_{2}+\be_{2}\otimes\be_{1})$, which
is trace free. Equation \eqr{10.24} gives
\[
\begin{aligned}
   \bT&=-p\I+\mu\gamma(\be_{1}\otimes\be_{2}+\be_{2}\otimes\be_{1}),
   \qquad\text{so}\qquad
   T_{12}\\
   &=\mu\gamma,\quad T_{11}=T_{22}=T_{33}=-p .
\end{aligned}
\]
The shear component is constitutive: measure it and you have measured
$\mu$. The three normal components are the reaction: measure them and you
have measured the boundary conditions.

\emph{Counting.} With $\bD=\Sym(\grad\bv)$, the linear momentum balance
is three equations, conservation of mass \eqr{9.84} is a fourth, and the
constraint $\dvg\bv=0$ of \eqr{9.85} is a fifth, for the five unknowns
$\rho$, $p$ and the three components of $\bv$. Remove $p$ and there are
five equations for four unknowns: an overdetermined, and generally
insoluble, system. The constraint added an equation, and the multiplier
it brought with it added exactly one unknown to pay for it. That balance
is not a coincidence of this example; it is what Lemma~\ref{lem:ch10ortho}
guarantees in general, one multiplier per independent constraint.
\end{example}

\begin{figure}[htbp]\centering
\renewcommand{\thefigure}{A}
\begin{tikzpicture}[scale=1.0,
   ar/.style={-{Stealth[length=2.0mm]},semithick}]
% ============ panel (a): open to the atmosphere ============
\begin{scope}
  \fill[blue!8] (0,0) rectangle (2.7,2.3);
  \draw[very thick] (0,2.7)--(0,0)--(2.7,0)--(2.7,2.7);
  \draw[thick,blue!55!black] (0,2.3)--(2.7,2.3);
  \node[font=\small,anchor=south] at (1.35,2.42) {$p=p_{a}$};
  \foreach \y/\L in {2.0/0.30, 1.55/0.52, 1.10/0.74, 0.65/0.96, 0.20/1.18}
     {\draw[ar,red!65!black] (\L,\y)--(0,\y);}
  \draw[ar] (3.05,0)--(3.05,2.3) node[right,font=\small] {$x_{3}$};
  \node[font=\small,anchor=north] at (1.35,-0.25) {(a)  $p(h)=p_{a}$};
\end{scope}
% ============ panel (b): sealed and pressed ============
\begin{scope}[xshift=6.0cm]
  \fill[blue!8] (0,0) rectangle (2.7,2.3);
  \draw[very thick] (0,2.7)--(0,0)--(2.7,0)--(2.7,2.7);
  \fill[gray!35] (0,2.3) rectangle (2.7,2.62);
  \draw[thick] (0,2.3) rectangle (2.7,2.62);
  \draw[ar,thick] (1.35,3.30)--(1.35,2.72);
  \node[font=\small,anchor=west] at (1.48,3.05) {$\Delta$};
  \foreach \y/\L in {2.0/0.72, 1.55/0.94, 1.10/1.16, 0.65/1.38, 0.20/1.60}
     {\draw[ar,red!65!black] (\L,\y)--(0,\y);}
  \node[font=\small,anchor=north] at (1.35,-0.25)
     {(b)  $p(h)=p_{a}+\Delta$};
\end{scope}
\end{tikzpicture}
\caption{The reaction pressure of Example~\ref{ex:ch10hydro}. In both
panels the fluid is the same, the motion is the same, namely rest, and
the constitutive part $\bG$ of the stress is the same, since $\bD=\Ob$.
The arrows show the pressure the fluid exerts on the left wall, growing
linearly with depth in both cases. Only the boundary condition at the top
differs, and the whole pressure field shifts by the constant $\Delta$
throughout the body. A constitutive equation is a function of the motion
and could not tell the two panels apart; the pressure is therefore not
constitutive, but a reaction fixed by the equations of motion and the
data.}
\label{fig:ch10press}
\end{figure}

\subsubsection*{(2) Inextensibility}

For inextensibility of embedded fibres with unit tangent $\bM$ in
$\kappa$ we use \eqr{10.12}. Since $\bM$ is a material vector,
$\dot\bM=\bze$ by Remark~\ref{rem:matvec}, so differentiating
$\ip\bM{\bC\bM}-1$ leaves only the term in which $\bC$ is
differentiated, and
\begin{equation}
   g_{\bC}\cdot\dot\bC=\ip\bM{\dot\bC\bM}=\bM\otimes\bM\cdot\dot\bC
   \tag{10.25}\label{eq:10.25}
\end{equation}
for all symmetric $\dot\bC$. The second equality is the component
identity
\[
   (\bM\otimes\bM)\cdot\bA=M_{A}M_{B}A_{AB}=M_{A}(\bA\bM)_{A}
   =\ip\bM{\bA\bM},
\]
using the component rule for the tensor product, Definition~\ref{def:tp},
and the component form of the tensor inner product. Thus, by the same
arbitrariness argument used for \eqr{10.21}, and noting that
$\bM\otimes\bM$ is symmetric,
\begin{equation}
   g_{\bC}=\bM\otimes\bM .
   \tag{10.26}\label{eq:10.26}
\end{equation}
Writing $\lambda=\sigma$, we use \eqr{10.18} to obtain
\begin{equation}
   \bN=\sigma\bF(\bM\otimes\bM)\bF^{t}=\sigma\,\bF\bM\otimes\bF\bM ,
   \tag{10.27}\label{eq:10.27}
\end{equation}
the second equality being the dyad rule $\bA(\ba\otimes\bb)\bB^{t}
=\bA\ba\otimes\bB\bb$, which follows from Lemma~\ref{lem:dyadprod}, or
directly in components:
$A_{iC}(M_{C}M_{D})B_{jD}=(A_{iC}M_{C})(B_{jD}M_{D})$. This reduces to
\begin{equation}
   \bN=\sigma\,\bm\otimes\bm
   \tag{10.28}\label{eq:10.28}
\end{equation}
when the constraint is in effect, where $\bm$ is the unit tangent to a
fibre in $\kappa_{t}$. The step is \eqr{2.67}, $\bF\bM=\mu\bm$, together
with the constraint $\mu=1$: with $\mu\neq1$ one would have
$\bF\bM\otimes\bF\bM=\mu^{2}\bm\otimes\bm$, and the factor $\mu^{2}$ is
absorbable into $\sigma$ in any case, which is why the constraint may be
imposed here without loss. Thus the reactive stress is mechanically
equivalent to a uniaxial stress along the fibres, exactly the state of
stress of \eqr{6.130}, with $\sigma$ the fibre tension.

Worklessness is again immediate, and this time it is Chapter 4 that
supplies it. From \eqr{10.28} and the identity used in \eqr{10.25},
\[
   \bN\cdot\bD=\sigma\,(\bm\otimes\bm)\cdot\bD=\sigma\,\ip\bm{\bD\bm}
   =\sigma\,(\ln\mu)^{\textstyle\cdot}=0 ,
\]
the third equality being \eqr{4.15} and the last the constraint
$\mu\equiv1$, whose logarithm is constant. A tension works only through
a stretching of the line it acts along, and the constraint is that there
is none. Compare the incompressible case: there the pair was pressure
against rate of volume change, here it is tension against rate of fibre
stretch, and in both the constraint annihilates the second member of the
pair.

\begin{example}[what an inextensible fibre still permits]
\label{ex:ch10fibre}
It is worth seeing that the constraint is far from paralysing. Take the
fibres along $\bM=\hbE_{1}$ and let the body undergo the simple shear of
Section~\ref{sec:shear} \emph{parallel} to them,
\[
   x_{1}=X_{1}+\gamma X_{2},\qquad x_{2}=X_{2},\qquad x_{3}=X_{3},
   \qquad\text{so}\qquad
   \bF=\shift+\gamma\,\be_{1}\otimes\hbE_{2}.
\]
Then $\bF\hbE_{1}=\shift\hbE_{1}+\gamma\be_{1}(\ip{\hbE_{2}}{\hbE_{1}})
=\be_{1}$, since $\ip{\hbE_{2}}{\hbE_{1}}=0$ and $\shift\hbE_{1}
=\be_{1}$ by \eqr{2.68}. Comparing with $\bF\bM=\mu\bm$ gives $\mu=1$ and
$\bm=\be_{1}$: the constraint holds for \emph{every} $\gamma$, however
large. Equivalently $C_{11}=\ip{\bF\hbE_{1}}{\bF\hbE_{1}}=1$.

Shear the same body across the fibres instead, $x_{2}=X_{2}+\gamma
X_{1}$, and $\bF\hbE_{1}=\be_{1}+\gamma\be_{2}$, so
$\mu=\sqrt{1+\gamma^{2}}\neq1$ for $\gamma\neq0$: forbidden. One
constraint, one direction lost; five of the six dimensions of $\Sym$
remain, which is Figure~\ref{fig:ch10ortho} in a concrete case.

The reaction in the permitted motion is $\bN=\sigma\be_{1}\otimes
\be_{1}$, with $\sigma(\bx,t)$ a free field. Its worklessness can be
checked without any general theory. Let the shear proceed in time,
$\gamma=\gamma(t)$. Then $\bv=\dot\gamma X_{2}\be_{1}
=\dot\gamma x_{2}\be_{1}$, since $x_{2}=X_{2}$, so
$\bL=\grad\bv=\dot\gamma\,\be_{1}\otimes\be_{2}$ and by \eqr{4.7}
\[
   \bD=\tfrac12\dot\gamma\big(\be_{1}\otimes\be_{2}
       +\be_{2}\otimes\be_{1}\big),
   \qquad\text{so}\qquad
   D_{11}=\ip{\be_{1}}{\bD\be_{1}}=0 ,
\]
and $\bN\cdot\bD=\sigma D_{11}=0$. What $\sigma$ actually is at each
point is settled, as always, by $\dvg\bT+\rho\bb=\rho\dot\bv$ and the
tractions applied to the ends of the fibres, not by any property of the
fibre.
\end{example}

\begin{figure}[htbp]\centering
\renewcommand{\thefigure}{B}
\begin{tikzpicture}[scale=1.0,
   ar/.style={-{Stealth[length=2.2mm]},thick}]
% ============ reference configuration ============
\begin{scope}
  \draw[thick] (0,0) rectangle (3,2);
  \foreach \y in {0.4,0.8,1.2,1.6}
     {\draw[gray!65] (0,\y)--(3,\y);}
  \draw[very thick,red!65!black] (0.9,1.2)--(1.9,1.2);
  \fill (0.9,1.2) circle (1.4pt);  \fill (1.9,1.2) circle (1.4pt);
  \node[font=\footnotesize,anchor=south,text=red!65!black]
     at (1.4,1.26) {length $1$};
  \draw[ar] (1.9,0.55)--(2.6,0.55) node[right=1pt,font=\small] {$\bM$};
  \node[font=\small,anchor=north] at (1.5,-0.2) {$\kappa$};
\end{scope}
% ============ current configuration: shear along the fibres ============
\begin{scope}[xshift=5.4cm]
  \draw[thick] (0,0)--(3,0)--(4,2)--(1,2)--cycle;
  \foreach \y/\a/\b in {0.4/0.2/3.2, 0.8/0.4/3.4, 1.2/0.6/3.6, 1.6/0.8/3.8}
     {\draw[gray!65] (\a,\y)--(\b,\y);}
  \draw[very thick,red!65!black] (1.5,1.2)--(2.5,1.2);
  \fill (1.5,1.2) circle (1.4pt);  \fill (2.5,1.2) circle (1.4pt);
  \node[font=\footnotesize,anchor=south,text=red!65!black]
     at (2.0,1.26) {length $1$};
  \draw[ar] (2.0,0.55)--(2.7,0.55) node[right=1pt,font=\small] {$\bm$};
  % the fibre tension: a cut normal to m, with opposed tractions
  \draw[guide] (2.90,0.02)--(2.90,1.98);
  \draw[ar,blue!55!black] (2.96,1.78)--(3.56,1.78);
  \draw[ar,blue!55!black] (2.84,1.78)--(2.24,1.78);
  \node[font=\small,text=blue!55!black,anchor=south] at (3.26,1.84)
     {$\sigma\bm$};
  \node[font=\small,anchor=north] at (2.0,-0.2) {$\kappa_{t}$};
\end{scope}
\end{tikzpicture}
\caption{An inextensible fibre family, $\bM$ in $\kappa$ and $\bm$ in
$\kappa_{t}$, under a shear parallel to the fibres. Every fibre keeps its
length, so the constraint $\mu=1$ holds for every amount of shear, while
the matrix between the fibres shears freely: the constraint removes one
dimension from $\Sym$ and leaves five. The reaction \eqr{10.28} is the
uniaxial stress $\sigma\bm\otimes\bm$, drawn as the pair of opposed
tractions on a cut normal to $\bm$. It does no work because
$\ip\bm{\bD\bm}=(\ln\mu)^{\textstyle\cdot}=0$, and its magnitude
$\sigma$ is not a property of the fibre but whatever the equations of
motion and the applied end tractions require.}
\label{fig:ch10fibre}
\end{figure}

\subsubsection*{(3) Rigidity}

In the case of rigidity it follows from \eqr{10.13} that, for each
ordered pair $(A,B)$,
\begin{equation}
\begin{aligned}
   (g_{AB})_{\bC}\cdot\dot\bC
   &=\dot C_{AB}=\big(\ip{\hbE_{A}}{\bC\hbE_{B}}\big)^{\textstyle\cdot}
    =\ip{\hbE_{A}}{\dot\bC\hbE_{B}}\\
   &=(\hbE_{A}\otimes\hbE_{B})\cdot\dot\bC
    =\Sym(\hbE_{A}\otimes\hbE_{B})\cdot\dot\bC ,
\end{aligned}
   \tag{10.29}\label{eq:10.29}
\end{equation}
on account of the symmetry of $\dot\bC$, where
\begin{equation}
   \Sym(\hbE_{A}\otimes\hbE_{B})
   =\tfrac12\big(\hbE_{A}\otimes\hbE_{B}+\hbE_{B}\otimes\hbE_{A}\big).
   \tag{10.30}\label{eq:10.30}
\end{equation}
The first line of \eqr{10.29} used $\dot{\dd}_{AB}=0$ and the fixity of
the referential basis, $\dot{\hbE}_{A}=\bze$; the second used the
identity established after \eqr{10.25}; the third dropped
$\Skw(\hbE_{A}\otimes\hbE_{B})$ by Lemma~\ref{lem:symskw}, since
$\dot\bC$ is symmetric. Thus
\begin{equation}
   (g_{AB})_{\bC}=\Sym(\hbE_{A}\otimes\hbE_{B})
   \tag{10.31}\label{eq:10.31}
\end{equation}
and it follows from \eqr{10.19} that
\begin{equation}
\begin{aligned}
   \bN&=\lambda_{AB}\bQ\big[\Sym(\hbE_{A}\otimes\hbE_{B})\big]\bQ^{t}
   =\bQ\big[\lambda_{(AB)}\hbE_{A}\otimes\hbE_{B}\big]\bQ^{t},\\
   &\text{with}\quad
   \lambda_{(AB)}=\tfrac12(\lambda_{AB}+\lambda_{BA}),
\end{aligned}
   \tag{10.32}\label{eq:10.32}
\end{equation}
where the usual summation convention is in effect, the constraint has
been imposed in the form $\bF=\bQ$, a rotation, and the $\lambda_{AB}$
are arbitrary scalars.

The passage from the first form of \eqr{10.32} to the second is a
relabelling worth doing once. Using \eqr{10.30},
\[
   \lambda_{AB}\Sym(\hbE_{A}\otimes\hbE_{B})
   =\tfrac12\lambda_{AB}\,\hbE_{A}\otimes\hbE_{B}
    +\tfrac12\lambda_{AB}\,\hbE_{B}\otimes\hbE_{A}.
\]
In the second sum both $A$ and $B$ are summed over, so they may be
interchanged as names, turning it into
$\tfrac12\lambda_{BA}\,\hbE_{A}\otimes\hbE_{B}$. Adding,
\[
   \lambda_{AB}\Sym(\hbE_{A}\otimes\hbE_{B})
   =\tfrac12(\lambda_{AB}+\lambda_{BA})\,\hbE_{A}\otimes\hbE_{B}
   =\lambda_{(AB)}\,\hbE_{A}\otimes\hbE_{B}.
\]
Only the symmetric part of the array $(\lambda_{AB})$ survives, which is
the same phenomenon as the symmetry of $g_{\bC}$: the skew part is
invisible because it is contracted against symmetric objects.

The term in square brackets to the right of the second equality in
\eqr{10.32} is an arbitrary symmetric referential tensor: it has six free
entries $\lambda_{(AB)}$, and every symmetric referential tensor is of
that form. Conjugation $\bS\mapsto\bQ\bS\bQ^{t}$ is linear, carries
symmetric tensors to symmetric tensors, since
$(\bQ\bS\bQ^{t})^{t}=\bQ\bS^{t}\bQ^{t}$, and is invertible with inverse
$\bT\mapsto\bQ^{t}\bT\bQ$. It therefore maps the symmetric referential
tensors \emph{onto} the symmetric spatial tensors. Thus $\bN$ is an
arbitrary symmetric spatial tensor. It follows that the net stress $\bT$
of \eqr{10.14} is constitutively indeterminate: the stress in a rigid
body is indeterminate. Nevertheless, the equations of rigid-body dynamics
of Section 6.4.2 suffice to determine the motion in terms of the applied
forces and torques, and vice versa.

\begin{remark}[why rigidity determines nothing, in one
line]\label{rem:ch10rigid}
The orthogonality condition \eqr{10.15} is the only thing that ever
restricted $\bN$, and in a rigid motion it says nothing at all. By
Remark~\ref{rem:Drigid}, $\bD=\Ob$ in a rigid motion, so
$\bN\cdot\bD=0$ holds for \emph{every} symmetric $\bN$ and the
requirement is vacuous. Equation \eqr{10.32} is that observation written
out. This also settles the count: six independent constraints remove all
six dimensions of admissible stretching and return all six dimensions of
$\Sym$ as reaction, which is the extreme case of the arithmetic in
Figure~\ref{fig:ch10ortho}.

The same fact appears in the energy balance. Chapter 6 observes that for
a rigid body the stress power is $\bT\cdot\bL=\bT\cdot\bW$, which
vanishes by the symmetry of $\bT$ and Lemma~\ref{lem:symskw}, so the
thermal and mechanical problems decouple. A material that can store no
strain energy can also tell you nothing about its stress.
\end{remark}

%=======================================================================
\setcounter{problemenv}{0}
\section{Problems}\label{sec:problems:c10}
%=======================================================================

%-----------------------------------------------------------------------
\begin{problem}[The three constraints as $O^{+}$ writes them]
Write the constraints of incompressibility, inextensibility and rigidity
in the frame of observer $O^{+}$.
\end{problem}

\begin{solution}
The general rule is \eqr{10.10}: whatever $\bar g_{\kappa}$ is, the
function used by $O^{+}$ is
\[
   \bar g^{+}_{\kappa^{+}}(\bC^{+})=\bar g_{\kappa}(\bK^{t}\bC^{+}\bK),
\]
and $\bC^{+}=\bK\bC\bK^{t}$ by \eqr{7.30}$_{1}$, with $\bK$ the fixed
orthogonal tensor of \eqr{7.15}. So each case is a substitution followed
by a cancellation of $\bK$'s.

\textbf{Incompressibility.} From \eqr{10.11},
\[
\begin{aligned}
   \bar g^{+}_{\kappa^{+}}(\bC^{+})&=\det(\bK^{t}\bC^{+}\bK)-1
   \\
   &=(\det\bK^{t})(\det\bC^{+})(\det\bK)-1
   =(\det\bK)^{2}\det\bC^{+}-1 ,
\end{aligned}
\]
and $(\det\bK)^{2}=1$ because $\bK$ is orthogonal, so
\[
   \boxed{\ \bar g^{+}_{\kappa^{+}}(\bC^{+})=\det\bC^{+}-1 .\ }
\]
The function has exactly the same form. Its value is also the same
number, since $\det\bC^{+}=(\det\bK)^{2}\det\bC=\det\bC$. Equivalently,
$J^{+}=1$, which is \eqr{7.19} with $J=1$.

\textbf{Inextensibility.} From \eqr{10.12}, moving $\bK$ across the inner
product by the definition of the transpose,
\[
   \bar g^{+}_{\kappa^{+}}(\bC^{+})
   =\ip\bM{(\bK^{t}\bC^{+}\bK)\bM}-1
   =\ip{\bK\bM}{\bC^{+}\bK\bM}-1 .
\]
By \eqr{7.36} the counterpart in $\hat{\mathbb{E}}^{3+}$ of the
referential unit vector $\bM$ is $\bM^{+}=\bK\bM$, so
\[
   \boxed{\ \bar g^{+}_{\kappa^{+}}(\bC^{+})
   =\ip{\bM^{+}}{\bC^{+}\bM^{+}}-1,\qquad \bM^{+}=\bK\bM .\ }
\]
Again the same form, with the fibre direction carried to the new
reference configuration by $\bK$. That $\bM^{+}$ is still a unit vector
is orthogonality of $\bK$: $|\bK\bM|^{2}=\ip{\bK\bM}{\bK\bM}
=\ip\bM{\bK^{t}\bK\bM}=\ip\bM{\hI\bM}=1$.

\textbf{Rigidity.} From \eqr{10.13} the constraint is $\bC=\hI$, so
\[
   \bC^{+}=\bK\hI\bK^{t}=\bK\bK^{t}=\hI^{+},
\]
which is \eqr{7.26}, the transformation law for the referential identity.
Componentwise, resolve on the basis $\hbE^{+}_{A}=\bK\hbE_{A}$, which is
orthonormal for the same reason $\bM^{+}$ was a unit vector. Then
\[
\begin{aligned}
   C^{+}_{AB}&=\ip{\hbE^{+}_{A}}{\bC^{+}\hbE^{+}_{B}}
   =\ip{\bK\hbE_{A}}{\bK\bC\bK^{t}\bK\hbE_{B}}
   \\
   &=\ip{\hbE_{A}}{\bK^{t}\bK\bC\hbE_{B}}
   =\ip{\hbE_{A}}{\bC\hbE_{B}}=C_{AB},
\end{aligned}
\]
so the components are literally the same numbers and
\[
   \boxed{\ \bar g^{+}_{(\kappa^{+})AB}=C^{+}_{AB}-\dd_{AB}.\ }
\]

\emph{What the three answers say together.} In each case the constraint
keeps its form and its numerical value; only the referential objects it
is written with, $\bM$ and $\hbE_{A}$, are carried over by $\bK$. That is
precisely the requirement \eqr{10.3} from which the whole of
Section~\ref{sec:ch10obs} was deduced, now verified in the three cases
that matter. Had any of the three come out differently, that constraint
would have failed the test of \eqr{10.9} and would have had to be
discarded, in the manner of Remark~\ref{rem:ch10rotblind}.
\end{solution}

%-----------------------------------------------------------------------
\begin{problem}[The reactive stress seen by another observer]
Find an expression for the reactive stress $\bN^{+}$ as perceived by
$O^{+}$, associated with a generic constraint. Show that all observers
agree on the values of the various constitutively indeterminate
multipliers.
\end{problem}

\begin{solution}
Everything follows from \eqr{10.18} once the three ingredients
$g^{+}_{\bC^{+}}$, $\bF^{+}$ and the transformation of $\bN$ itself are
in hand.

\textbf{Step 1: the gradient transforms by $\bK$.} We claim
\[
   g^{+}_{\bC^{+}}=\bK\,g_{\bC}\,\bK^{t}.
\]
By definition, $g^{+}_{\bC^{+}}$ is the tensor for which
$\dl g^{+}=g^{+}_{\bC^{+}}\cdot\dl\bC^{+}$ for every symmetric increment
$\dl\bC^{+}$. From \eqr{10.10}, $g^{+}(\bC^{+})=g(\bK^{t}\bC^{+}\bK)$, so
the increments are related by $\dl\bC=\bK^{t}\,\dl\bC^{+}\bK$, $\bK$
being fixed. Hence, for any symmetric $\bA$,
\[
\begin{aligned}
   \dl g^{+}&=g_{\bC}\cdot(\bK^{t}\bA\bK)
   =\tr\big(g_{\bC}(\bK^{t}\bA\bK)^{t}\big)
   \\
   &=\tr\big(g_{\bC}\bK^{t}\bA^{t}\bK\big)
   =\tr\big(\bK g_{\bC}\bK^{t}\bA^{t}\big)
   =(\bK g_{\bC}\bK^{t})\cdot\bA ,
\end{aligned}
\]
where the fourth equality moved the trailing $\bK$ to the front. Both
traces are defined, the first on $\hVt$ and the second on
$\hat{\mathbb{E}}^{3+}$, and they are equal for the same component
reason checked after \eqr{10.17}:
the two products are the same sum of products of the same numbers.
Comparing the two expressions for $\dl g^{+}$ and using the uniqueness of
the gradient gives the claim. Note that $\bK g_{\bC}\bK^{t}$ is symmetric
whenever $g_{\bC}$ is, as it must be.

\textbf{Step 2: the sandwich is unchanged up to $\bQ$.} With
$\bF^{+}=\bQ\bF\bK^{t}$ from \eqr{7.17},
\[
\begin{aligned}
   \bF^{+}g^{+}_{\bC^{+}}(\bF^{+})^{t}
   &=(\bQ\bF\bK^{t})(\bK g_{\bC}\bK^{t})(\bQ\bF\bK^{t})^{t}
   =\bQ\bF\bK^{t}\bK\,g_{\bC}\,\bK^{t}\bK\bF^{t}\bQ^{t}\\
   &=\bQ\bF\hI\,g_{\bC}\,\hI\bF^{t}\bQ^{t}
   =\bQ\big(\bF g_{\bC}\bF^{t}\big)\bQ^{t}.
\end{aligned}
\]
Every $\bK$ has cancelled, which had to happen: the choice of reference
configuration is bookkeeping and cannot survive in a spatial tensor. So
by \eqr{10.18} as written by $O^{+}$,
\[
   \boxed{\ \bN^{+}=\lambda^{+}\bF^{+}g^{+}_{\bC^{+}}(\bF^{+})^{t}
   =\lambda^{+}\,\bQ\big(\bF g_{\bC}\bF^{t}\big)\bQ^{t}.\ }
\]

\textbf{Step 3: $\bN$ transforms as a spatial tensor.} By \eqr{7.40} the
net stress obeys $\bT^{+}=\bQ\bT\bQ^{t}$, and by \eqr{8.35} the
constitutive part obeys $\bG^{+}=\bQ\bG\bQ^{t}$, which is exactly the
statement that the two observers agree about material response. Since
$\bN=\bT-\bG$ and $\bN^{+}=\bT^{+}-\bG^{+}$,
\[
   \bN^{+}=\bQ\bT\bQ^{t}-\bQ\bG\bQ^{t}=\bQ(\bT-\bG)\bQ^{t}
   =\bQ\bN\bQ^{t}
   =\lambda\,\bQ\big(\bF g_{\bC}\bF^{t}\big)\bQ^{t}.
\]

\textbf{Step 4: the multipliers agree.} Subtracting the two expressions
for $\bN^{+}$,
\[
   (\lambda^{+}-\lambda)\,\bQ\big(\bF g_{\bC}\bF^{t}\big)\bQ^{t}=\Ob .
\]
Now $\bF g_{\bC}\bF^{t}\neq\Ob$, since $\bF$ is invertible and
$g_{\bC}\neq\hOb$ for a non-degenerate constraint, and $\bQ$ is
invertible, so the conjugated tensor is non-zero. A scalar multiple of a
non-zero tensor vanishes only if the scalar does. Therefore
\[
   \boxed{\ \lambda^{+}=\lambda .\ }
\]

\textbf{Several constraints.} With $n$ constraints, Steps 1 and 2 apply
to each $g^{(i)}$ separately, so \eqr{10.19} as written by $O^{+}$ reads
$\bN^{+}=\sum_{i}\lambda^{+}_{i}\bQ(\bF g^{(i)}_{\bC}\bF^{t})\bQ^{t}$,
while Step 3 gives $\bN^{+}=\sum_{i}\lambda_{i}\bQ(\bF
g^{(i)}_{\bC}\bF^{t})\bQ^{t}$. Subtracting,
\[
   \sum_{i=1}^{n}(\lambda^{+}_{i}-\lambda_{i})\,
   \bQ\big(\bF g^{(i)}_{\bC}\bF^{t}\big)\bQ^{t}=\Ob ,
\]
and since conjugation by $\bQ$ is an invertible linear map it preserves
linear independence, so the tensors
$\bQ(\bF g^{(i)}_{\bC}\bF^{t})\bQ^{t}$ are linearly independent whenever
the constraints are independent. Every coefficient must therefore vanish:
$\lambda^{+}_{i}=\lambda_{i}$ for each $i$.

\emph{The three cases.} For incompressibility $\lambda=-\bar p$, so
$\bar p^{+}=\bar p$. The thermodynamic pressure is also observer
independent, since $\tilde p=\rho^{2}\psi_{\rho}$ by \eqr{9.74} and both
$\rho$ and $\psi$ are absolute scalars by \eqr{7.37}; hence
$p^{+}=\tilde p^{+}+\bar p^{+}=p$, which is Problem 4 of Chapter 9. For
inextensibility $\sigma^{+}=\sigma$: the fibre tension is a number both
observers read off the same. For rigidity $\lambda^{+}_{AB}
=\lambda_{AB}$, although in that case the statement is empty of content,
since by Remark~\ref{rem:ch10rigid} the multipliers are arbitrary to
begin with.

\emph{Why this had to come out this way.} A multiplier is a scalar, and
Chapter 7 assumed absolute scalars to be the same for all observers. The
calculation above is not an independent discovery of that principle but a
consistency check on \eqr{10.18}: had $\lambda^{+}$ differed from
$\lambda$, the constraint reaction would have been observer dependent
while the constraint itself was not, and \eqr{10.14} would have split the
stress differently for different observers.
\end{solution}

%-----------------------------------------------------------------------
\begin{problem}[A laminate of inextensible sheets]
Consider a laminated material formed from thin sheets of stiff paper
interspersed in a soft matrix. Take the sheets to be continuously
distributed parallel planes spanned by $\hbE_{1}$ and $\hbE_{2}$ in the
reference configuration, so that the constraints are
$\ip{\hbE_{1}}{\bC\hbE_{2}}=0$ and
$\ip{\hbE_{1}}{\bC\hbE_{1}}=\ip{\hbE_{2}}{\bC\hbE_{2}}=1$. Show that
there can be no extensional or shear strain in the plane of the sheets,
but that transverse normal and shear strains are permitted. What is the
form of the reactive stress?
\end{problem}

\begin{solution}
The three constraints are statements about the upper left two by two
block of the matrix $(C_{AB})$:
\[
   C_{11}=C_{22}=1,\qquad C_{12}=C_{21}=0 ,
\]
where $C_{AB}=\ip{\hbE_{A}}{\bC\hbE_{B}}$ as in Chapter 2. Nothing is
said about $C_{13}$, $C_{23}$ or $C_{33}$.

\textbf{No in-plane strain, in any direction.} Take any material fibre
lying in a sheet, $\bM=M_{1}\hbE_{1}+M_{2}\hbE_{2}$ with
$M_{1}^{2}+M_{2}^{2}=1$. Its stretch follows from \eqr{2.83},
\[
\begin{aligned}
   \mu^{2}&=\ip\bM{\bC\bM}=M_{A}M_{B}C_{AB}
   \\
   &=M_{1}^{2}C_{11}+2M_{1}M_{2}C_{12}+M_{2}^{2}C_{22}
   =M_{1}^{2}+M_{2}^{2}=1 ,
\end{aligned}
\]
so $\mu=1$ for every in-plane direction: not just the two coordinate
directions but every fibre in the sheet is inextensible. For the angles,
\eqr{2.105} gives the cosine of the angle between the images of two
reference fibres $\bM$, $\bM'$ as
$\ip\bM{\bC\bM'}/(\mu\mu')$. With $\bM=\hbE_{1}$ and $\bM'=\hbE_{2}$ this
is $C_{12}/(1\cdot1)=0$: the right angle between the two coordinate
fibres is preserved, so there is no in-plane shear either. In terms of
the Lagrange strain $\bE=\tfrac12(\bC-\hI)$ of \eqr{2.97} the three
statements are simply
\[
   E_{11}=E_{22}=E_{12}=0 .
\]
Since $\bE$ restricted to the sheet plane vanishes identically, the sheet
deforms rigidly in its own plane. Its local area is preserved as well:
with $\ba_{\alpha}=\bF\hbE_{\alpha}$ for $\alpha=1,2$ we have
$\ip{\ba_{1}}{\ba_{1}}=C_{11}=1$, $\ip{\ba_{2}}{\ba_{2}}=C_{22}=1$ and
$\ip{\ba_{1}}{\ba_{2}}=C_{12}=0$, so $\{\ba_{1},\ba_{2}\}$ is an
orthonormal pair and $|\ba_{1}\times\ba_{2}|=1$.

\textbf{What is still permitted.} The entries $C_{13}$, $C_{23}$,
$C_{33}$ are unconstrained, so $E_{13}$, $E_{23}$, $E_{33}$ are free:
the sheets may slide over one another, tilt relative to their common
normal, and move apart or together. Three of the six components of $\bE$
are frozen and three are free, which is the count of
Remark~\ref{rem:ch10sixmax} with $n=3$.

\textbf{The reactive stress.} Write the three constraint functions as
\[
   g^{(1)}=\ip{\hbE_{1}}{\bC\hbE_{1}}-1,\quad
   g^{(2)}=\ip{\hbE_{2}}{\bC\hbE_{2}}-1,\quad
   g^{(3)}=\ip{\hbE_{1}}{\bC\hbE_{2}} .
\]
Their gradients follow from the computation that produced \eqr{10.26}
and \eqr{10.31}, since $g^{(1)}$ and $g^{(2)}$ are instances of
\eqr{10.12} and $g^{(3)}$ is an instance of \eqr{10.13} with $A\neq B$:
\[
   g^{(1)}_{\bC}=\hbE_{1}\otimes\hbE_{1},\qquad
   g^{(2)}_{\bC}=\hbE_{2}\otimes\hbE_{2},\qquad
   g^{(3)}_{\bC}=\Sym(\hbE_{1}\otimes\hbE_{2}) .
\]
Substituting in \eqr{10.19} and using the dyad rule
$\bF(\ba\otimes\bb)\bF^{t}=\bF\ba\otimes\bF\bb$ from \eqr{10.27},
\[
   \bN=\lambda_{1}\ba_{1}\otimes\ba_{1}+\lambda_{2}\ba_{2}\otimes\ba_{2}
      +\tfrac12\lambda_{3}\big(\ba_{1}\otimes\ba_{2}
      +\ba_{2}\otimes\ba_{1}\big).
\]
Because $\{\ba_{1},\ba_{2}\}$ was shown above to be orthonormal, write
$\bm_{1}=\ba_{1}$, $\bm_{2}=\ba_{2}$ and let $\bm_{3}=\bm_{1}\times
\bm_{2}$ be the unit normal to the deformed sheet, so that
$\{\bm_{1},\bm_{2},\bm_{3}\}$ is an orthonormal basis of $\Vt$. Then
\[
   \boxed{\ \bN=\lambda_{1}\bm_{1}\otimes\bm_{1}
   +\lambda_{2}\bm_{2}\otimes\bm_{2}
   +\tfrac12\lambda_{3}\big(\bm_{1}\otimes\bm_{2}
   +\bm_{2}\otimes\bm_{1}\big).\ }
\]
Equivalently: $\bN$ is an arbitrary symmetric tensor with
$\bN\bm_{3}=\bze$, that is, an arbitrary state of \emph{plane stress in
the plane of the deformed sheets}. It has three free components, one per
constraint.

\emph{Check that it does no work.} By \eqr{10.17} the admissible
stretchings are those with $\ip{\bm_{1}}{\bD\bm_{1}}
=\ip{\bm_{2}}{\bD\bm_{2}}=\ip{\bm_{1}}{\bD\bm_{2}}=0$, that is, those
whose in-plane block vanishes. The reaction has only an in-plane block.
Resolving both on $\{\bm_{i}\}$, $\bN\cdot\bD=N_{ij}D_{ij}$ and every
term has at least one factor from a vanishing block, so $\bN\cdot\bD=0$.
The reaction lives exactly where the stretching is forbidden, and the
stretching survives exactly where the reaction is absent. That
complementarity is Figure~\ref{fig:ch10ortho} again, now with a
three-dimensional space of reactions and a three-dimensional space of
admissible stretchings inside the six-dimensional $\Sym$.

\emph{Reading the answer.} The reaction is the membrane stress carried by
the stiff paper: two normal components in the sheet and one in-plane
shear. The matrix is free to shear transversely and to change thickness,
and the reaction supplies no stress for those modes, so the transverse
response of the laminate is entirely constitutive, which is why a stack
of paper is stiff in its plane and floppy across it.
\end{solution}

%-----------------------------------------------------------------------
\begin{problem}[A laminate whose sheets preserve area]
The body is laminated as in the previous problem, but now constrained in
such a way that the planes experience no change in local surface area as
they deform. Find the form of the reactive stress.
\end{problem}

\begin{solution}
The single constraint has to be expressed through $\bC$ before anything
else can be done, and Chapter 2 has the tool.

\textbf{Step 1: the constraint.} A material surface element of the sheet
has reference unit normal $\hbE_{3}$, since the sheets are spanned by
$\hbE_{1}$ and $\hbE_{2}$. By the Piola--Nanson formula \eqr{2.54},
$\bn\,\dl a=\bF^{*}\hbE_{3}\,\dl A$, so taking magnitudes of both sides
and using $|\bn|=1$,
\[
   \frac{\dl a}{\dl A}=\big|\bF^{*}\hbE_{3}\big|
   =\sqrt{\ip{\hbE_{3}}{(\bF^{*})^{t}\bF^{*}\hbE_{3}}} .
\]
The tensor in the middle is the cofactor of $\bC$. Indeed, by
Proposition~\ref{prop:I2cof}, $\bF^{*}=J\bF^{-t}$, so
\[
   (\bF^{*})^{t}\bF^{*}=\big(J\bF^{-t}\big)^{t}\big(J\bF^{-t}\big)
   =J^{2}\bF^{-1}\bF^{-t}=J^{2}(\bF^{t}\bF)^{-1}
   =(\det\bC)\bC^{-1}=\bC^{*},
\]
using $\bF^{-1}\bF^{-t}=(\bF^{t}\bF)^{-1}$ from
Proposition~\ref{prop:prodinv}, $J^{2}=\det\bC$ from \eqr{10.11}, and
$\bC^{*}=(\det\bC)\bC^{-1}$ from the derivation of \eqr{10.21}. This is
Remark~\ref{rem:cofgeom} in tensor form: the cofactor is the tensor that
carries area elements. Hence $\dl a=\dl A$ is $\ip{\hbE_{3}}{\bC^{*}
\hbE_{3}}=1$, and
\[
   g(\bC)=\ip{\hbE_{3}}{\bC^{*}\hbE_{3}}-1=C^{*}_{33}-1
   =C_{11}C_{22}-C_{12}^{2}-1 ,
\]
the last step because the $(3,3)$ entry of the cofactor matrix is the
determinant of the complementary two by two block, and $C_{21}=C_{12}$.
So the constraint is $\det\bC_{\parallel}=1$, where $\bC_{\parallel}$
denotes the in-plane block. Note in passing that the three constraints of
Problem 3 imply this one, since they give $\det\bC_{\parallel}
=1\cdot1-0=1$; the present constraint is the weaker requirement that the
in-plane block have unit determinant rather than be the identity.

\textbf{Step 2: the gradient.} Let $g_{\bC}=G_{AB}\hbE_{A}\otimes
\hbE_{B}$ with $G$ symmetric. For a symmetric increment $\dl\bC$,
\[
   g_{\bC}\cdot\dl\bC=G_{11}\dl C_{11}+G_{22}\dl C_{22}+G_{33}\dl C_{33}
   +2G_{12}\dl C_{12}+2G_{13}\dl C_{13}+2G_{23}\dl C_{23},
\]
the factors of two arising because each off-diagonal pair contributes
twice to the sum $G_{AB}\dl C_{AB}$. Differentiating
$g=C_{11}C_{22}-C_{12}^{2}-1$ directly,
\[
   \dl g=C_{22}\,\dl C_{11}+C_{11}\,\dl C_{22}-2C_{12}\,\dl C_{12}.
\]
Matching the coefficient of each independent increment,
\[
   G_{11}=C_{22},\quad G_{22}=C_{11},\quad G_{12}=-C_{12},\quad
   G_{33}=G_{13}=G_{23}=0 ,
\]
that is,
\[
   g_{\bC}=C_{22}\,\hbE_{1}\otimes\hbE_{1}+C_{11}\,\hbE_{2}\otimes\hbE_{2}
   -C_{12}\big(\hbE_{1}\otimes\hbE_{2}+\hbE_{2}\otimes\hbE_{1}\big).
\]
The array $(G_{AB})$ is the adjugate of the in-plane block, extended by
zeros, so
\[
   g_{\bC}=\big(\det\bC_{\parallel}\big)\bC_{\parallel}^{-1}
   \quad\text{extended by zero},
\]
which is \eqr{10.21} one dimension down, as one would expect from a
constraint that is a determinant in two dimensions rather than three. On
the constraint set $\det\bC_{\parallel}=1$ this reduces to
$\bC_{\parallel}^{-1}$, the analogue of \eqr{10.22}. As always the
reduction is performed only after the differentiation, per
Remark~\ref{rem:ch10apost}.

\textbf{Step 3: the reaction.} With $\ba_{\alpha}=\bF\hbE_{\alpha}$,
$\alpha=1,2$, and the dyad rule of \eqr{10.27},
\[
   \bF g_{\bC}\bF^{t}
   =\big(\bC_{\parallel}^{-1}\big)_{\alpha\beta}\,
     \ba_{\alpha}\otimes\ba_{\beta},
\]
the Greek indices running over $1,2$ only. Now observe that
$\bC_{\parallel}$ is precisely the Gram matrix of the pair
$\{\ba_{1},\ba_{2}\}$, since
$\ip{\ba_{\alpha}}{\ba_{\beta}}=\ip{\bF\hbE_{\alpha}}{\bF\hbE_{\beta}}
=\ip{\hbE_{\alpha}}{\bC\hbE_{\beta}}=C_{\alpha\beta}$. The tensor
$\bP=(\bC_{\parallel}^{-1})_{\alpha\beta}\ba_{\alpha}\otimes\ba_{\beta}$
is then the orthogonal projection onto the plane they span. To verify
this, apply it to $\ba_{\gamma}$:
\[
   \bP\ba_{\gamma}
   =\big(\bC_{\parallel}^{-1}\big)_{\alpha\beta}
     \ba_{\alpha}\,\ip{\ba_{\beta}}{\ba_{\gamma}}
   =\big(\bC_{\parallel}^{-1}\big)_{\alpha\beta}
     C_{\beta\gamma}\,\ba_{\alpha}
   =\dd_{\alpha\gamma}\ba_{\alpha}=\ba_{\gamma},
\]
and to any $\bw$ perpendicular to both:
$\bP\bw=(\bC_{\parallel}^{-1})_{\alpha\beta}\ba_{\alpha}
\ip{\ba_{\beta}}\bw=\bze$. A symmetric tensor that fixes a plane and
annihilates its orthogonal complement is the orthogonal projection onto
that plane, by Definition~\ref{def:splits} and the uniqueness of the
splitting. Writing $\bm_{3}$ for the unit normal to the deformed sheet,
$\bP=\I-\bm_{3}\otimes\bm_{3}$, and \eqr{10.18} gives
\[
   \boxed{\ \bN=\lambda\big(\I-\bm_{3}\otimes\bm_{3}\big).\ }
\]

\emph{Reading the answer.} Compare \eqr{10.23}. Incompressibility is the
statement that a volume element does not change, and its reaction is an
isotropic pressure in three dimensions, $-\bar p\I$. The present
constraint is the statement that an area element does not change, and its
reaction is an isotropic pressure in the two dimensions of the sheet,
$\lambda(\I-\bm_{3}\otimes\bm_{3})$. One constraint, one multiplier, and
the reaction is the identity of the space in which the measure is
frozen.

\emph{Consistency with Problem 3.} The reaction found there had three
free components in the sheet plane; the reaction found here has one, and
it is the isotropic member of that three-parameter family, obtained by
setting $\lambda_{1}=\lambda_{2}=\lambda$ and $\lambda_{3}=0$. That is
the right relationship: Problem 3 imposes three constraints and Problem 4
imposes one of the consequences of those three, so it must buy a
one-dimensional subspace of the earlier three-dimensional space of
reactions. Weaker constraint, smaller reaction, and the arithmetic of
Remark~\ref{rem:ch10sixmax} is respected in both.

\emph{Worklessness, checked.} Admissible stretchings satisfy
$\bP\cdot\bD=0$ by \eqr{10.17}, and $\bP\cdot\bD
=\tr(\bP\bD)=D_{11}+D_{22}$ resolved on $\{\bm_{i}\}$. So the constraint
permits any $\bD$ whose in-plane trace vanishes, which is the statement
that the sheet may stretch in one in-plane direction provided it
contracts in the other, keeping the area fixed. And
$\bN\cdot\bD=\lambda\bP\cdot\bD=0$ for exactly those $\bD$. A
two-dimensional pressure works only through a change of area, and the
constraint is that there is none.
\end{solution}
